% chapters/09/section-06.tex
\section{Yang Belum Ada pada Mesin Keadaan Ini}
\label{sec:belum-ada-fsm}

NPC pada bab ini sudah memiliki empat keadaan, tetapi masih belum dapat
mengelilingi rintangan. Bagian ini menjelaskan batas kemampuannya dan apa yang
ditambahkan pada bab berikutnya.

\subsection{Keputusan tanpa Pengetahuan tentang Jalan}

Mesin keadaan menjawab pertanyaan ``apa yang sebaiknya saya lakukan?'', bukan
pertanyaan ``bagaimana cara menuju ke sana?''. NPC pada bab ini memutuskan
untuk mengejar, dan pelaksanaannya hanya satu, yaitu bergerak lurus ke arah
pemain.

Akibatnya, kelemahan yang sudah Anda amati pada Bab~VIII masih ada. Empat
keadaan tidak menolong NPC yang terhalang dinding: keadaan mengejar tetap
menghasilkan gerak lurus ke arah pemain, sehingga NPC tetap menekan dinding
selama pemain berada di baliknya. Yang bertambah hanyalah keragaman perilakunya,
bukan keterangan tentang jalan.

Pekerjaan itu menjadi bagian Bab~X. Di sana, keadaan mengejar tidak lagi
menghasilkan gerak lurus, melainkan memanggil pencarian jalur untuk menentukan
rute menuju pemain.

\begin{catatabox}
Perhatikan pemisahan tanggung jawab yang mulai terbentuk. Mesin keadaan
menentukan \textbf{niat}, sedangkan pencarian jalur menentukan \textbf{rute}.
Pemisahan itu persis seperti yang dijelaskan pada Bagian~2.5 sebagai lapisan
keputusan dan lapisan navigasi. Pada Bab~X kedua lapisan itu akhirnya
berhubungan.
\end{catatabox}

\subsection{Bentuk Mesin Keadaan yang Lebih Besar}

Mesin keadaan pada bab ini cukup untuk permainan sederhana. Pada permainan yang
lebih besar, dua bentuk pengembangan berikut lazim dipakai.

\begin{itemize}
  \item \textbf{Mesin keadaan bertingkat} (\istilah{hierarchical state
        machine}). Sebuah keadaan dapat memuat mesin keadaan yang lebih kecil di
        dalamnya. Misalnya, keadaan ``mengejar'' dapat memuat keadaan
        ``menghindari rintangan'' dan ``mendekati pemain''. Cara itu mencegah
        jumlah keadaan bertambah secara berlipat.
  \item \textbf{Pohon perilaku} (\istilah{behaviour tree}). Perilaku disusun
        sebagai pohon keputusan yang dibaca dari atas ke bawah, dan setiap
        cabangnya menyatakan satu pertimbangan. Cara itu lebih mudah dibaca
        ketika jumlah perilakunya banyak, dan banyak dipakai pada permainan
        yang sesungguhnya. Contoh penerapannya dapat dilihat pada \mbox{[R26]}.
\end{itemize}

Kedua bentuk itu berada di luar cakupan mata kuliah ini. Yang perlu Anda ketahui
adalah bahwa keduanya bertolak dari gagasan yang sama dengan mesin keadaan:
memisahkan keadaan yang mungkin terjadi dari syarat perpindahannya.

\subsection{Ringkasan Perkembangan}

\begin{table}[htbp]
\centering
\caption{Bagian yang ditambahkan pada perilaku NPC pada bab-bab berikutnya}
\label{tab:perkembangan-npc}
{\setlength{\tabcolsep}{4pt}\footnotesize%
\begin{tabularx}{\textwidth}{@{}C{1.6cm} L{5.0cm} Y@{}}
\toprule
\headrow \textbf{Bab} & \textbf{Yang ditambahkan} & \textbf{Akibatnya} \\
\midrule
Bab~X & Pencarian jalur A* pada peta berhalangan &
Keadaan mengejar tidak lagi menghasilkan gerak lurus, melainkan rute yang
menghindari rintangan. \\
\addlinespace
Bab~XI & Penggabungan mesin keadaan dan pencarian jalur &
NPC mengejar dengan rute memutar, dan biaya perhitungannya diatur agar tetap
ringan. \\
\addlinespace
Bab~XIII & Fisika mesin pada Unity &
Sebagian perilaku dapat disusun memakai perkakas mesin, sedangkan logika
keputusannya tetap Anda tulis. \\
\bottomrule
\end{tabularx}}
\end{table}

Dengan gambaran tersebut, tampak bahwa bab ini melengkapi bagian pertama
lapisan keputusan yang telah diperkenalkan pada Bab~VIII. Setelah bab ini, Anda
memiliki seluruh perkakas untuk menyusun perilaku NPC yang beragam. Yang masih
kurang hanyalah kemampuan NPC untuk mengetahui jalannya sendiri.
